\chapter{Gastrointestinal Bleeding}

\section*{Definition and Overview}
Gastrointestinal (GI) bleeding is bleeding anywhere in the digestive tract, from the mouth and oesophagus through the stomach and intestines to the anus. It is divided into upper GI bleeding, from the oesophagus, stomach, or first part of the small intestine, and lower GI bleeding, from the lower small intestine, colon, or rectum. GI bleeding can range from a slow, silent loss of blood that causes anaemia to a sudden, massive haemorrhage that is life-threatening. The common causes include peptic ulcers, gastritis, and oesophageal varices (upper), and haemorrhoids, diverticular disease, polyps, inflammatory bowel disease, and cancer (lower). Any GI bleeding warrants medical assessment, and significant bleeding is an emergency. This chapter explains the causes, symptoms, and management.

\section*{Epidemiology}
Gastrointestinal bleeding is a common emergency. Upper GI bleeding affects around 50--100 people per 100,000 per year, and it is the most common reason for emergency gastroscopy; peptic ulcers and gastritis, often from NSAIDs or H.\ pylori, cause most cases, and variceal bleeding is a major risk in liver disease. Lower GI bleeding is also common, with diverticular disease, haemorrhoids, and bowel cancer among the causes. While most bleeding is minor, the overall mortality of significant GI bleeding is around 5--10%, largely in older people, which is why it is taken seriously.

\section*{Aetiology and Pathophysiology}
\begin{itemize}
    \item \textbf{Upper GI bleeding:} peptic ulcers (from H.\ pylori, NSAIDs, and stress), gastritis and erosions, oesophagitis, Mallory--Weiss tears (from violent vomiting), oesophageal varices (in liver cirrhosis), and stomach cancer
    \item \textbf{Lower GI bleeding:} haemorrhoids and anal fissures, diverticular disease, angiodysplasia, polyps and bowel cancer, inflammatory bowel disease, and infectious colitis
    \item \textbf{Slow bleeding:} chronic, low-volume blood loss from ulcers, polyps, or cancer causes iron-deficiency anaemia without dramatic symptoms
    \item \textbf{Acute bleeding:} sudden, large-volume loss causes symptoms of shock, including faintness, pallor, and collapse
    \item \textbf{Blood in the stool:} upper bleeding produces black, tarry stools (melena); lower bleeding produces bright or dark red blood
    \item \textbf{Vomiting blood:} haematemesis, which may be fresh blood or "coffee-ground" material
    \item \textbf{The danger:} rapid blood loss threatens circulation and organ perfusion; variceal bleeding is particularly dangerous
    \item \textbf{Anticoagulants and antiplatelets:} medicines such as warfarin, apixaban, and aspirin increase the risk and severity of bleeding
\end{itemize}

\section*{Clinical Features}
\begin{itemize}
    \item Vomiting fresh blood or dark, coffee-ground material
    \item Black, tarry, foul-smelling stools (melena) from upper bleeding
    \item Bright red blood or clots from the back passage from lower bleeding
    \item Blood mixed with the stool, or blood on the toilet paper
    \item Faintness, dizziness, pallor, sweating, and a rapid heart rate with significant loss
    \item Fatigue, breathlessness, and pale skin from chronic anaemia
    \item Abdominal pain, depending on the cause
    \item Collapse and shock in severe haemorrhage
    \item The absence of visible blood does not exclude bleeding; slow loss causes anaemia
\end{itemize}

\section*{Differential Diagnosis}
The source of bleeding is identified by its pattern and tests.
\begin{itemize}
    \item Upper versus lower GI bleeding, distinguished by the appearance of the stool and vomiting
    \item Peptic ulcers, gastritis, oesophagitis, and varices (upper)
    \item Haemorrhoids, fissures, diverticular disease, polyps, inflammatory bowel disease, and cancer (lower)
    \item Bleeding from other sources, such as a bleeding disorder or medicines
    \item False alarms: red foods and iron tablets can darken stool, and beetroot can colour it red
    \item Endoscopy and imaging locate the bleeding source
\end{itemize}

\section*{When to Seek Medical Care}
Any GI bleeding, however minor it seems, should be assessed by a doctor, because the amount passed does not reflect the total lost and the cause needs diagnosis. Significant bleeding is an emergency.

\textbf{Call 000 (or local emergency services) for:}
\begin{itemize}
    \item Vomiting blood or coffee-ground material
    \item Black, tarry stools
    \item Large amounts of bright red blood from the back passage
    \item Faintness, dizziness, pallor, or collapse
    \item A rapid heart rate or breathlessness with bleeding
    \item Bleeding in a person taking blood thinners
\end{itemize}

\section*{Diagnosis and Investigations}
\begin{itemize}
    \item \textbf{Resuscitation first:} in significant bleeding, fluids and blood products are given while the cause is investigated
    \item \textbf{Blood tests:} full blood count, blood group and cross-match, clotting tests, kidney and liver function
    \item \textbf{Gastroscopy:} the key test for upper bleeding, allowing diagnosis and treatment, such as injecting or clipping a bleeding ulcer
    \item \textbf{Colonoscopy:} for lower bleeding, with treatment of polyps and bleeding sites
    \item \textbf{CT angiography:} for bleeding too fast for endoscopy to see clearly
    \item \textbf{Capsule endoscopy:} a swallowed camera for obscure small bowel bleeding
    \item \textbf{Anaemia work-up:} iron studies when slow bleeding is suspected
\end{itemize}

\section*{Management}
\begin{itemize}
    \item \textbf{Emergency care:} intravenous fluids, blood transfusion, and monitoring in significant bleeding
    \item \textbf{Endoscopic treatment:} injection, clips, or heat to stop bleeding ulcers; banding of varices; removal of polyps
    \item \textbf{Acid suppression:} intravenous PPIs for upper bleeding
    \item \textbf{Variceal treatment:} medicines, endoscopic banding, and specialist care for liver disease
    \item \textbf{Surgery or radiology:} for bleeding that does not stop with endoscopy
    \item \textbf{Treat the cause:} eradicate H.\ pylori, stop NSAIDs, and treat the underlying condition
    \item \textbf{Manage medicines:} review anticoagulants and antiplatelets with a doctor
    \item \textbf{Treat anaemia:} iron replacement after bleeding has stopped
    \item \textbf{Follow-up:} repeat endoscopy and surveillance as planned
\end{itemize}

\section*{Complications}
\begin{itemize}
    \item Shock and organ failure from rapid blood loss
    \item Death, mainly in older people and those with liver disease
    \item Iron-deficiency anaemia from chronic bleeding
    \item Re-bleeding, especially from ulcers and varices
    \item Complications of endoscopy and surgery
    \item Progression of the underlying cause, such as bowel cancer, if not diagnosed
\end{itemize}

\section*{Prognosis and Outcomes}
The outcome of GI bleeding has improved greatly with modern endoscopy. Most upper GI bleeding is stopped by endoscopic treatment, and the overall mortality, while significant in emergencies, is largely in older and high-risk patients. Chronic bleeding is diagnosed and treated, and iron deficiency is corrected. Bowel cancer detected through bleeding and investigated early is often curable. The keys to a good outcome are treating bleeding as an emergency, stopping the cause, and managing risk factors.

\section*{Prevention and Self-Care}
\begin{itemize}
    \item Use NSAIDs at the lowest dose for the shortest time, and take stomach protection if on long-term treatment
    \item Treat H.\ pylori infection
    \item Do not smoke and limit alcohol
    \item Do the National Bowel Cancer Screening test
    \item Report any blood in the stool or black stools to a doctor promptly
    \item Have unexplained anaemia investigated
    \item Discuss anticoagulant medicines with a doctor, and never stop them without advice
    \item Seek urgent care for significant bleeding
\end{itemize}

\section*{Sources and Further Reading}
The following reputable sources provide further information for healthcare students and patients seeking reliable, current advice about gastrointestinal bleeding.
\begin{itemize}
    \item Healthdirect Australia. \textit{Gastrointestinal bleeding.} https://www.healthdirect.gov.au/gastrointestinal-bleeding
    \item The Gut Foundation. \textit{Bleeding from the gut.} https://www.gutfoundation.org.au/
    \item Gastroenterology Society of Australia (GESA). \textit{Endoscopy and GI bleeding.} https://www.gesa.org.au/
    \item NHS. \textit{Gastrointestinal bleeding.} https://www.nhs.uk/
    \item Mayo Clinic. \textit{Gastrointestinal bleeding.} https://www.mayoclinic.org/
    \item MSD Manual. \textit{Gastrointestinal bleeding.} https://www.msdmanuals.com/
\end{itemize}
